\section{结论与展望}\label{sec:conclusion}

% --- 5.1 主要发现 ---

\subsection{主要发现}

本文系统综述了金属有机框架材料在光催化水解产氢领域的最新研究进展，主要发现总结如下。

第一，MOF材料凭借其高比表面积、可调能带结构和多孔晶态骨架，已成为光催化水解产氢领域的重要研究平台。
自2010年首次报道以来，MOF光催化产氢研究取得了显著进展~\citep{chen2023mof_owsp_review}。 %% CLAIM_ID: EC-2026-001
不同金属节点（Zr、Ti、Fe等）对光吸收性质和催化活性具有显著影响，为材料设计提供了丰富的调控空间~\citep{reddy2021mof_h2_advances}。 %% CLAIM_ID: EC-2026-019

第二，异质结构建是提升MOF光催化性能的核心策略。
Type-II、Z-scheme和S-scheme三种异质结机制各有特点：
Type-II型有效延长载流子寿命，Z-scheme和S-scheme则在保持强氧化还原电位方面更具优势~\citep{zhang2025optimization_strategies}。 %% CLAIM_ID: EC-2026-006
MOF-on-MOF异质外延结构为全分解水提供了新的实现路径~\citep{uio66_mil88b_2024}。 %% CLAIM_ID: EC-2026-004

第三，配体工程取得了突破性进展。
三配体Ti-MOF通过多配体协同优化实现了产氢速率的大幅提升~\citep{liu2024three_ligand_timof}。 %% CLAIM_ID: EC-2026-020
供体-受体分子工程利用分子内电荷转移机制，使产氢活性提升了数个数量级~\citep{yan2023donor_acceptor_mof}。 %% CLAIM_ID: EC-2026-021
这些进展表明，从分子层面精确调控MOF的电子结构是实现高效光催化的关键路径。

第四，MOF衍生材料兼具前驱体的结构优势和无机材料的高稳定性，在产氢半反应中展现出优异性能~\citep{singh2025mini_review}。 %% CLAIM_ID: EC-2026-017
配体功能化、金属掺杂和等离子体耦合等策略有效拓展了MOF的光吸收范围和载流子分离效率~\citep{li2024mof_composites_h2}。 %% CLAIM_ID: EC-2026-008

第五，在代表性实验成果方面，CFA-Zn全分解水体系（AQE $3.09\%$，稳定性$>100~\mathrm{h}$）和RuO$_2$/N,S-TiO$_2$异质结（AQY $10.0\%$）代表了当前领域的较高水平~\citep{cfazn2025ows,rtta2024heterojunction}。 %% CLAIM_ID: EC-2026-003

% --- 5.2 当前挑战 ---

\subsection{当前面临的挑战}

尽管MOF光催化水解产氢取得了长足进步，该领域仍面临若干关键挑战。

\textbf{（1）水稳定性和光稳定性。}
部分MOF材料（尤其是含Zn、Cu节点的MOF）在水溶液中的结构稳定性不足，长时间光照可导致骨架坍塌和活性衰减~\citep{hassan2024emerging_trends}。 %% CLAIM_ID: EC-2026-013
提升MOF在反应条件下的长期稳定性是实现实际应用的先决条件。

\textbf{（2）量子效率瓶颈。}
当前MOF光催化全分解水的AQE普遍低于$5\%$，与理论极限和实际应用需求之间存在较大差距~\citep{chen2023mof_owsp_review}。 %% CLAIM_ID: EC-2026-002
载流子在MOF骨架中的复合损失和表面反应动力学缓慢是主要制约因素。

\textbf{（3）规模化合成与成本。}
MOF材料的合成通常依赖溶剂热法，产量有限、成本较高，且部分有机配体价格昂贵~\citep{wang2025mof_pec_review}。 %% CLAIM_ID: EC-2026-005
\citet{jiang2026mof_review}强调，发展绿色、低成本的MOF规模化合成方法是推动产业化的关键。 %% CLAIM_ID: EC-2026-018

\textbf{（4）机理理解不足。}
MOF光催化水分解的微观机制——包括光生载流子在有机-无机杂化骨架中的传输路径、表面活性位点的精确识别以及水分子在孔道中的吸附-反应动力学——仍需深入研究~\citep{pmc2024mof_energy_env}。

% --- 5.3 未来发展方向 ---

\subsection{未来发展方向}

基于上述分析，本文提出以下未来研究方向。

\textbf{全光谱响应设计。}
通过上转换材料耦合、窄带隙组分复合和光热协同催化等策略，
将MOF光催化体系的响应范围从紫外-可见光区拓展至近红外区域，
充分利用太阳光谱中$95\%$以上的能量~\citep{mof_fullspectrum_2025}。 %% CLAIM_ID: EC-2026-011

\textbf{微孔限域激子转移工程。}
基于微孔尺寸与激子扩散长度匹配的研究发现~\citep{zhou2023micropore_exciton}，
有望通过精确调控MOF孔道几何构型来优化激子传输路径，
抑制非辐射复合损失，实现光催化效率的本质提升。 %% CLAIM_ID: EC-2026-022

\textbf{产氢耦合高值化学品合成。}
将产氢半反应与有机物选择性氧化反应耦合，
在热力学上降低全反应过电位，在经济性上提升附加值，
有望推动MOF光催化从基础研究走向实际应用~\citep{lin2025coupled_oxidation}。 %% CLAIM_ID: EC-2026-012

\textbf{MOF-MXene等新型复合体系。}
MXene材料优异的电导率和亲水性与MOF的高孔隙率形成互补，
MOF-MXene复合体系在载流子传输和光吸收方面展现出协同增强效应~\citep{mxene_mof_2024}。 %% CLAIM_ID: EC-2026-014
类似的跨材料平台整合策略值得进一步探索。

\textbf{供体-受体分子工程深化。}
D-A型MOF和三配体策略的成功表明，从分子层面精确调控电荷转移路径具有巨大潜力~\citep{yan2023donor_acceptor_mof,liu2024three_ligand_timof}。 %% CLAIM_ID: EC-2026-021
结合计算化学预筛选和高通量实验验证，有望加速发现高性能MOF光催化剂。

\textbf{绿色合成与产业化路径。}
发展水相合成、机械化学合成和连续流合成等绿色方法，
降低MOF的生产成本和环境负荷，
为大规模制氢应用奠定基础~\citep{jiang2026mof_review}。 %% CLAIM_ID: EC-2026-018

\textbf{原位表征与理论计算。}
利用原位光谱（时间分辨光致发光、瞬态吸收光谱）和密度泛函理论（DFT）计算，
深入解析MOF光催化水分解的微观机制，
建立构效关系模型，指导高性能催化剂的理性设计。 %% NO_CITE: common knowledge
